% 康奈尔听课笔记模板（配色取自王道课件）
\documentclass[10.5pt,a4paper,oneside]{article}
\usepackage[a4paper,top=2.3cm,bottom=2.2cm,left=5.7cm,right=1.5cm,
  marginparwidth=4.25cm,marginparsep=0.55cm,headheight=14pt,footskip=1.1cm]{geometry}
\usepackage{xeCJK}
\usepackage{xeCJKfntef}
\xeCJKDeclareCharClass{FullLeft}{"201C,"2018}
\xeCJKDeclareCharClass{FullRight}{"201D,"2019}
\xeCJKDeclareCharClass{CJK}{"2460->"24FF,"2605,"2606}
\setmainfont{Noto Sans CJK SC}[BoldFont={Noto Sans CJK SC Bold},ItalicFont={Noto Sans CJK SC},ItalicFeatures={FakeSlant=0.18}]
\setsansfont{Noto Sans CJK SC}[BoldFont={Noto Sans CJK SC Bold}]
\setmonofont{DejaVu Sans Mono}[Scale=0.9]
\setCJKmainfont{Noto Sans CJK SC}[BoldFont={Noto Sans CJK SC Bold},ItalicFont={Noto Sans CJK SC},ItalicFeatures={FakeSlant=0.18}]
\setCJKsansfont{Noto Sans CJK SC}[BoldFont={Noto Sans CJK SC Bold}]
\setCJKmonofont{Noto Sans Mono CJK SC}[Scale=0.9]
\linespread{1.32}
\usepackage{amsmath,amssymb}
\usepackage[dvipsnames,table]{xcolor}
\usepackage{graphicx}
\usepackage{adjustbox}
\usepackage{tikz}
\usetikzlibrary{calc}
\usepackage[most]{tcolorbox}
\usepackage{eso-pic}
\usepackage{fancyhdr}
\usepackage{titlesec}
\usepackage{needspace}
\usepackage{enumitem}
\usepackage{xltabular,booktabs,array,colortbl}
\usepackage{marginnote}
\reversemarginpar
\renewcommand*{\raggedleftmarginnote}{}
\renewcommand*{\raggedrightmarginnote}{}
\usepackage{caption}
\usepackage{titletoc}
\usepackage{fancyvrb}
\usepackage{framed}
\usepackage[section]{placeins}
\usepackage{hyperref}
\hypersetup{colorlinks=true,linkcolor=wdblue,urlcolor=wdblue,pdftitle={1.1\_1
计算机网络的概念},pdfauthor={听课笔记}}

% ---------- 课件配色 ----------
\definecolor{wdblue}{HTML}{4472C4}
\definecolor{wdlightblue}{HTML}{5B9BD5}
\definecolor{wdtitle}{HTML}{4A86C8}
\definecolor{wddark}{HTML}{1F3864}
\definecolor{wdorange}{HTML}{ED7D31}
\definecolor{wdgray}{HTML}{A5A5A5}
\definecolor{wdgold}{HTML}{FFC000}
\definecolor{wdgreen}{HTML}{70AD47}
\definecolor{wddeepgreen}{HTML}{3F7F2A}
\definecolor{wdred}{HTML}{C00000}
\definecolor{wdpurple}{HTML}{7030A0}
\definecolor{cuebg}{HTML}{F2F6FC}
\definecolor{codebg}{HTML}{F7F9FC}
\color[HTML]{262626}

% ---------- 行内样式 ----------
\newcommand{\hlred}[1]{\textcolor{wdred}{\bfseries #1}}
\newcommand{\hlblue}[1]{\textcolor{wdblue}{\bfseries #1}}
\newcommand{\hlgreen}[1]{\textcolor{wddeepgreen}{\bfseries #1}}
\newcommand{\hlorange}[1]{\textcolor{wdorange}{\bfseries #1}}
\newcommand{\hlpurple}[1]{\textcolor{wdpurple}{\bfseries #1}}
\newcommand{\uln}[1]{\CJKunderline[format=\color{wdred},thickness=0.9pt,depth=0.28em]{#1}}
\newcommand{\wavy}[1]{\CJKunderwave[format=\color{wdorange}]{#1}}
\newcommand{\hlmark}[1]{\CJKunderline[format=\color{wdgold!55},thickness=0.55em,depth=-0.05em]{#1}}
\newcommand{\tagbox}[2]{\tcbox[on line,colback=#1,colframe=#1,coltext=white,boxsep=0.5pt,left=3pt,right=3pt,top=0.5pt,bottom=0.5pt,arc=2pt,fontupper=\small\bfseries]{#2}}

% ---------- 康奈尔版式背景：左侧线索栏 ----------
\newlength{\cuestripw}\setlength{\cuestripw}{5.15cm}
\AddToShipoutPictureBG{%
  \begin{tikzpicture}[remember picture,overlay]
    \fill[cuebg] (current page.north west) rectangle ($(current page.south west)+(\cuestripw,0)$);
    \draw[wdblue,line width=1.2pt] ($(current page.north west)+(\cuestripw,-1.55cm)$) -- ($(current page.south west)+(\cuestripw,1.55cm)$);
    \node[anchor=west,text=wdblue,font=\small\bfseries] at ($(current page.north west)+(0.75cm,-1.25cm)$) {线索 · 问题};
    \node[anchor=west,text=wdblue,font=\small\bfseries] at ($(current page.north west)+(\cuestripw+0.55cm,-1.25cm)$) {笔记 · 老师主线 + 补充讲解};
  \end{tikzpicture}}

% ---------- 页眉页脚 ----------
\pagestyle{fancy}\fancyhf{}
\renewcommand{\headrulewidth}{0pt}
\fancyhead[R]{\small\color{wdgray}第1章 计算机网络体系结构}
\fancyfoot[C]{\small\color{wdgray}\thepage}
\fancyfoot[R]{\small\color{wdorange}1.1\_1 计算机网络的概念 · 听课笔记}

\fancypagestyle{plain}{\fancyhf{}\fancyfoot[C]{\small\color{wdgray}\thepage}\fancyfoot[R]{\small\color{wdorange}1.1\_1
计算机网络的概念 · 听课笔记}}

% ---------- 标题：仿课件的蓝色标题 + 六色条 ----------
\newcommand{\sixbar}{\par\nointerlineskip\vspace{3pt}\centering
  \begin{tikzpicture}\foreach \c [count=\i] in {wdlightblue,wdorange,wdgray,wdgold,wdblue,wdgreen}
     {\fill[\c] ({(\i-1)*1.05},0) rectangle ({(\i-1)*1.05+0.92},0.11);}\end{tikzpicture}\par}
\titleformat{\section}[block]{\centering\Large\bfseries\color{wdtitle}}{}{0pt}{}[\sixbar]
\titlespacing*{\section}{0pt}{16pt}{10pt}
\pretocmd{\section}{\Needspace{8\baselineskip}}{}{}
\titleformat{\subsection}{\Needspace{7\baselineskip}\large\bfseries\color{wddark}}{\textcolor{wdblue}{\rule[-0.15em]{4pt}{1.05em}}\hspace{6pt}}{0pt}{}
\titlespacing*{\subsection}{0pt}{12pt}{5pt}
\titleformat{\subsubsection}{\Needspace{5\baselineskip}\normalsize\bfseries\color{wdblue}}{\textcolor{wdorange}{\large\ensuremath{\blacktriangleright}}\hspace{4pt}}{0pt}{}
\titlespacing*{\subsubsection}{0pt}{8pt}{3pt}
\setcounter{secnumdepth}{0}
\setlength{\parindent}{0pt}\setlength{\parskip}{5pt}
\setlist{nosep,leftmargin=1.5em,itemsep=2pt}
\setlist[itemize,1]{label=\textcolor{wdblue}{\small\ensuremath{\bullet}}}
\setlist[itemize,2]{label=\textcolor{wdorange}{\scriptsize\ensuremath{\blacktriangleright}}}
\setlist[enumerate,1]{label=\textcolor{wdblue}{\bfseries\arabic*.}}
\providecommand{\tightlist}{\setlength{\itemsep}{2pt}\setlength{\parskip}{0pt}}

% ---------- 目录 ----------
\setcounter{tocdepth}{3}
\titlecontents{section}[0pt]{\addvspace{5pt}\bfseries\color{wddark}}{}{}{\titlerule*[0.7pc]{\color{wdgray!60}.}\contentspage}
\titlecontents{subsection}[1.6em]{\addvspace{1pt}\small\color{black!75}}{}{}{\titlerule*[0.7pc]{\color{wdgray!40}.}\contentspage}
\newcommand{\tocback}{\par\vspace{-3pt}{\hfill\footnotesize\hyperlink{toc}{\textcolor{wdgray}{\ensuremath{\uparrow}\,回到目录}}}\par\vspace{-4pt}}

% ---------- 线索栏 ----------
\newcommand{\cueblock}[2]{\Needspace{4\baselineskip}\leavevmode#1{#2}\par\vspace{-\baselineskip}\vspace{-\parskip}}
% 线索直接贴在段首行的左侧（不用 marginnote，避免分页时偶尔落进正文栏）
\newcommand{\cuebox}[1]{\leavevmode\llap{\smash{\begin{minipage}[t]{\marginparwidth}\raggedright\small\linespread{1.18}\selectfont #1\end{minipage}}\hspace{\marginparsep}}}
\newcommand{\cuenote}[1]{\cuebox{{\color{wdblue}\bfseries Q\,}\color{wddark}#1}}
\newcommand{\cuekey}[1]{\cuebox{\color{wdred}\bfseries #1}}

% ---------- 讲解框 ----------
\tcbset{wdbox/.style={enhanced,breakable,arc=2pt,boxrule=0pt,leftrule=3pt,left=7pt,right=7pt,top=5pt,bottom=5pt,
   before skip=7pt,after skip=7pt,fonttitle=\bfseries\small,attach title to upper={\quad},coltitle=#1,
   colframe=#1,colback=#1!6}}
\newtcolorbox{defbox}[1][]{enhanced,breakable,arc=3pt,boxrule=0.8pt,colframe=wdblue,colback=wdblue!4,
   left=8pt,right=8pt,top=6pt,bottom=6pt,fonttitle=\bfseries\small,coltitle=white,colbacktitle=wdblue,
   title={定义#1},before skip=8pt,after skip=8pt}
\newtcolorbox{thmbox}[1]{enhanced,breakable,arc=3pt,boxrule=0.8pt,colframe=wdblue,colback=wdblue!4,
   left=8pt,right=8pt,top=6pt,bottom=6pt,fonttitle=\bfseries\small,coltitle=white,colbacktitle=wdblue,
   title={#1},before skip=8pt,after skip=8pt}
\newtcolorbox{bridgebox}{wdbox=wdblue,title={承上启下}}
\newtcolorbox{plainbox}{wdbox=wddeepgreen,colframe=wdgreen,colback=wdgreen!8,title={大白话}}
\newtcolorbox{deepbox}{wdbox=wdpurple,title={背后的思想}}
\newtcolorbox{exambox}{wdbox=wdred,title={考点}}
\newtcolorbox{pitfallbox}{wdbox=wdorange,title={易错提醒}}
\newtcolorbox{bookbox}{wdbox=wddark,colback=wdgray!10,colframe=wdgray,title={王道书补充}}
\newtcolorbox{teacherbox}{wdbox=wdorange,colback=wdgold!8,colframe=wdgold,coltitle=wdorange,title={老师原话}}
\newtcolorbox{fixbox}{wdbox=wdred,colback=wdred!4,title={勘误 / 说明}}
% 习题笔记用
\newtcolorbox{questionbox}[1][]{enhanced,breakable,arc=3pt,boxrule=0.8pt,colframe=wddark,colback=white,
   left=8pt,right=8pt,top=6pt,bottom=6pt,fonttitle=\bfseries\small,coltitle=white,colbacktitle=wddark,
   title={#1},before skip=10pt,after skip=6pt}
\newtcolorbox{videobox}{wdbox=wdorange,colback=wdgold!8,colframe=wdgold,coltitle=wdorange,title={视频讲解}}
\newtcolorbox{jiexibox}{wdbox=wddark,colback=wdgray!10,colframe=wdgray,title={书上解析}}
\newtcolorbox{addbox}{wdbox=wdpurple,title={我的补充}}
\newtcolorbox{recallbox}{wdbox=wdblue,title={回顾正课}}
\newtcolorbox{watchbox}{enhanced,breakable,arc=3pt,boxrule=1pt,colframe=wdorange,colback=wdgold!12,
   fonttitle=\bfseries\small,coltitle=white,colbacktitle=wdorange,title={建议去看视频},left=8pt,right=8pt,before skip=8pt,after skip=8pt}
\newtcolorbox{summarybox}{enhanced,arc=3pt,
   boxrule=1pt,colframe=wdblue,colback=white,colbacktitle=wdblue,coltitle=white,fonttitle=\bfseries,
   title={总结区 · 用自己的话复述这一节},left=10pt,right=10pt,top=8pt,bottom=8pt,before skip=14pt}

% ---------- 代码 ----------

\makeatletter
\@ifundefined{Shaded}{}{\renewenvironment{Shaded}{\def\FrameCommand{\fboxsep=5pt\fboxrule=0.5pt\fcolorbox{wdblue!35}{codebg}}%
   \MakeFramed{\advance\hsize-\width\FrameRestore}\small}{\endMakeFramed}}
\makeatother
\fvset{fontsize=\small}

% ---------- 表格 ----------
\renewcommand{\arraystretch}{1.42}
\newcolumntype{Y}{>{\raggedright\arraybackslash}X}

% ---------- 图片 ----------
% 课件截图做成浮动体：LaTeX 会自动用正文填满页面，避免大块空白
\setcounter{topnumber}{3}\setcounter{bottomnumber}{2}\setcounter{totalnumber}{4}
\renewcommand{\topfraction}{0.85}\renewcommand{\bottomfraction}{0.6}\renewcommand{\textfraction}{0.12}\renewcommand{\floatpagefraction}{0.75}
\setlength{\textfloatsep}{10pt plus 3pt minus 2pt}\setlength{\floatsep}{8pt plus 2pt}
\newcommand{\slidefig}[3]{\begin{figure}[!ht]\centering\begin{tcolorbox}[enhanced,width=#2\linewidth,arc=2pt,boxrule=0.6pt,colframe=wdgray!70,
   colback=white,left=2pt,right=2pt,top=2pt,bottom=2pt,drop fuzzy shadow=black!25]\includegraphics[width=\linewidth]{#1}\end{tcolorbox}
   \vspace{-3pt}{\small\color{wdgray}#3}\end{figure}}
\newcommand{\plainfig}[3]{\par\noindent\begin{minipage}{\linewidth}\centering\includegraphics[width=#2\linewidth]{#1}\par\vspace{2pt}{\small\color{wdgray}#3}\end{minipage}\par\vspace{4pt}}

\begin{document}
% ---------- 封面条（仿课件封面） ----------
\thispagestyle{plain}
\begin{tcolorbox}[enhanced,enlarge left by=-4.8cm,width=\linewidth+4.8cm,arc=0pt,boxrule=0pt,
   colback=wddark,interior style={left color=wddark,right color=wdblue},left=18pt,right=18pt,top=14pt,bottom=12pt,
   overlay={\draw[white,line width=2.5pt] ([xshift=10pt,yshift=-10pt]frame.north west) -- ++(0,-0.9cm);
            \draw[white,line width=2.5pt] ([xshift=10pt,yshift=-10pt]frame.north west) -- ++(1.6cm,0);
            \draw[white,line width=2.5pt] ([xshift=-10pt,yshift=10pt]frame.south east) -- ++(0,0.9cm);
            \draw[white,line width=2.5pt] ([xshift=-10pt,yshift=10pt]frame.south east) -- ++(-1.6cm,0);}]
  {\color{white!85!wdlightblue}\small 计算机网络 · 第1章
计算机网络体系结构 · 康奈尔听课笔记}\par\vspace{6pt}
  {\color{white}\fontsize{22}{28}\selectfont\bfseries 1.1\_1
计算机网络的概念}\par\vspace{6pt}
  {\color{white!80!wdlightblue}\small 课件 1.1\_1 · 正课字幕 P2 · 王道书
§1.1.1（第 1 页）}
\end{tcolorbox}
\vspace{4pt}
\hypertarget{ux8fd9ux4e00ux7ae0ux5728ux5b66ux4ec0ux4e48}{%
\section{这一章在学什么}\label{ux8fd9ux4e00ux7ae0ux5728ux5b66ux4ec0ux4e48}}

\Needspace{4\baselineskip}

\cuenote{第 1 章在考试里考什么？第一次学觉得模糊怎么办？}第 1
章「计算机网络体系结构」会认识很多和这门课相关的概念，并\uln{建立整门课程的知识框架}。考试中这一章主要考\hlred{选择题和简单的计算题}。

很多概念第一次学会觉得难理解，这很正常：它们基本都会在后续章节深入学习。学完整门课再回来看这一章，会觉得很简单。

这一个视频解决两个问题：\textbf{什么是计算机网络}；\textbf{计算机网络、互连网、互联网}三个术语的区别（后两个读音相同，“连”和“联”不一样）。

\hypertarget{ux4e00ux4ec0ux4e48ux662fux8ba1ux7b97ux673aux7f51ux7edc}{%
\section{一、什么是计算机网络}\label{ux4e00ux4ec0ux4e48ux662fux8ba1ux7b97ux673aux7f51ux7edc}}

\cueblock{\cuenote}{\relax{计算机网络的定义里有哪几个关键词？“自治”是什么意思？}}

\begin{defbox}

\textbf{计算机网络}（computer
networking）：一个将众多\hlred{分散的}、\hlred{自治的}计算机系统，通过\uln{通信设备与线路}连接起来，由功能完善的\uln{软件}实现\hlred{资源共享和信息传递}的系统。

\end{defbox}

对照生活经验逐个拆开：

\begin{itemize}
\tightlist
\item
  \textbf{计算机系统}：手机、笔记本电脑、台式机、智能手表、智能电视都算；
\item
  \textbf{分散、自治}：这些设备彼此独立，\uln{其中一台坏掉不影响另一台的运行}；
\item
  \textbf{通信设备与线路}：家用路由器、5G
  基站是通信设备；线路可以是无线（WiFi、5G），也可以是有线（网线）；
\item
  \textbf{软件}：百度网盘、微信等，用它们共享资源、聊天、传文件，也就是\textbf{资源共享和信息传递}。
\end{itemize}

\slidefig{assets/1.1_1-p04.jpg}{0.6}{课件：什么是计算机网络}

\begin{bookbox}

王道书第 1
页给出的就是这条定义，并说明“一般认为”：计算机网络的定义没有统一标准，不同教材说法不同。

\end{bookbox}

\hypertarget{ux4e8cux8ba1ux7b97ux673aux7f51ux7edcux8282ux70b9-ux94feux8def}{%
\section{二、计算机网络：节点 +
链路}\label{ux4e8cux8ba1ux7b97ux673aux7f51ux7edcux8282ux70b9-ux94feux8def}}

\cueblock{\cuenote}{\relax{最简单的计算机网络长什么样？集线器、交换机各起什么作用？}}

\begin{defbox}

\textbf{计算机网络}（简称\textbf{网络}）：由若干\textbf{节点}（node）和连接这些节点的\textbf{链路}（link）组成。节点可以是计算机、集线器、交换机、路由器等；链路可以是有线链路、无线链路。

\end{defbox}

\begin{addbox}

\textbf{名词统一}：课件 1.1\_1
写“结点”，王道书写“节点（Node，或译为结点）“，课件 1.1\_5
也写”节点”。本笔记统一用\textbf{节点}。

\end{addbox}

老师用宿舍联机打《帝国时代》的经历，把网络从小搭到大：

\begin{enumerate}
\def\labelenumi{\arabic{enumi}.}
\tightlist
\item
  \textbf{两台电脑 + 一根网线}：2 个节点、1
  条链路，就是一个最简单的计算机网络；
\item
  \textbf{集线器}（Hub）：每台笔记本只有一个网口，多台电脑要联机，就用网线把
  4 台电脑都接到一个集线器上。这个网络有 \uln{5 个节点}（4 台电脑 + 1
  个集线器）、4
  条有线链路。缺点是集线器\hlred{不能避免数据冲突}：两台电脑同时发数据会“打架”，结果谁都发不出去，所以现在已很少用；
\item
  \textbf{交换机}（Switch）：长得和集线器很像，也能把多个节点连成一个网络，但\hlred{不会发生数据冲突}。现在家庭、公司、学校组建内部网络通常用交换机；交换机还能连交换机，组成校园网这样较大的内部网络。
\end{enumerate}

一个计算机网络通常画成一朵\textbf{云}。

\slidefig{assets/1.1_1-p07.jpg}{0.48}{课件：用集线器组建的简单网络}

\slidefig{assets/1.1_1-p08.jpg}{0.48}{课件：用交换机组建的简单网络}

\begin{teacherbox}

集线器的原理在第 2 章学（课件标注：工作在物理层），交换机在第 3
章学（工作在数据链路层），路由器在第 4
章学（工作在网络层）。现在只需要知道它们\textbf{能干什么}。

\end{teacherbox}

\hypertarget{ux4e09ux4e92ux8fdeux7f51internetux8defux7531ux5668ux628aux591aux4e2aux7f51ux7edcux8fdeux8d77ux6765}{%
\section{三、互连网（internet）：路由器把多个网络连起来}\label{ux4e09ux4e92ux8fdeux7f51internetux8defux7531ux5668ux628aux591aux4e2aux7f51ux7edcux8fdeux8d77ux6765}}

\cueblock{\cuenote}{\relax{不同的计算机网络之间怎么互联互通？银行内网为什么在家里访问不了？}}

\begin{defbox}

\textbf{路由器}（router）：可以把两个或多个计算机网络互相连接起来，形成规模更大的计算机网络，也可以称为\hlred{\textbf{互连网}（internet，“连接”的连）}。路由器还能连接别的路由器，让互连网的规模不断扩大。

\end{defbox}

\slidefig{assets/1.1_1-p10.jpg}{0.55}{课件：路由器把多个计算机网络连成互连网}

\textbf{例子：某银行的内部网络}。银行有总行和 A、B、C、D
四个分行，每个分行是一个独立的计算机网络，分行通过路由器和总行的数据中心互联互通。客户在
A 分行存 1000 元，存款信息经路由器逐层传到总行数据中心，下次在 D
分行取钱也能查到余额。这就是一个典型的互连网。

银行数据和钱直接挂钩，所以银行会建一个\uln{专属于自己的内部网络}；你家连的网络和银行内网没有互联互通，所以在家里访问不了。

\hypertarget{ux56dbux5bb6ux7528ux8defux7531ux5668-ux8bfeux7a0bux91ccux7684ux8defux7531ux5668}{%
\section{四、家用路由器 ≠
课程里的路由器}\label{ux56dbux5bb6ux7528ux8defux7531ux5668-ux8bfeux7a0bux91ccux7684ux8defux7531ux5668}}

\cueblock{\cuenote}{\relax{家用路由器可以看作哪两种设备的结合？手机投屏到电视，为什么要连同一个网络？}}

\begin{pitfallbox}

这门课里的“路由器”和你家的“家用路由器”有区别：\hlred{家用路由器 = 路由器 + 交换机 + 其他功能}。路由器这个词之后会频繁出现，不要把两者等同。

\end{pitfallbox}

\begin{itemize}
\tightlist
\item
  \textbf{交换机部分}：电视、电脑用网线接家用路由器，手机用 WiFi
  连上，逻辑上都相当于连在\uln{同一个交换机}上，组成你家内部的计算机网络。手机和电视处于同一个网络之下，所以可以\textbf{投屏}；
\item
  \textbf{路由器部分}：把家庭内部网络和外面的网络连接起来，于是家里的设备就能上网。路由器后面是一个网络，路由器外面是另一个网络。
\end{itemize}

\slidefig{assets/1.1_1-p12.jpg}{0.55}{课件：家庭内部网络（交换机组网，路由器连接外部网络）}

一句话对照：\hlred{交换机把多个\textbf{节点}连成一个网络；路由器把多个\textbf{网络}连在一起}。

\hypertarget{ux4e94ux4e92ux8054ux7f51internetux56e0ux7279ux7f51}{%
\section{五、互联网（Internet，因特网）}\label{ux4e94ux4e92ux8054ux7f51internetux56e0ux7279ux7f51}}

\Needspace{4\baselineskip}

\cuenote{ISP 是什么？互连网和互联网在通信协议上有什么区别？}既然路由器能把多个网络连起来，就可以把全世界的计算机网络都连起来。做这件事的是\hlred{\textbf{互联网服务提供商 ISP}（Internet Service Provider）}，如中国电信、中国移动、中国联通。全球很多这样的企业和国际机构，共同构建了覆盖全球的互联网，也就是\textbf{因特网}。中文里一般用“互联网”（联合的联）特指因特网。

\begin{itemize}
\tightlist
\item
  互联网对每个公众开放：\uln{只要交钱、交网费，就可以接入}。开通宽带那天，师傅把一根网线接到你家路由器上，另一头最终连到运营商的某个商用路由器上（中间其实还有更复杂的连接，课件做了简化），你家的网络就和整个互联网连起来了；
\item
  所以平时说的“\textbf{上网}”，这个“网”指的是互联网（Internet）。
\end{itemize}

\slidefig{assets/1.1_1-p14.jpg}{0.6}{课件：互联网（因特网，Internet）}

\begin{defbox}

\textbf{通信协议的要求}：

\begin{itemize}
\tightlist
\item
  \hlred{互联网（Internet）必须使用 \textbf{TCP/IP 协议族}作为通信规则}；
\item
  互连网（internet）可以使用\uln{任意通信协议}。
\end{itemize}

\end{defbox}

为什么？老师用秦始皇“书同文”作比：统一六国前各国文字不同，秦国人只能给秦国人写信，相当于各国的“通信协议”各不相同；统一文字后，咸阳的人和邯郸的人就能按同一规则互相通信。覆盖全世界的互联网要让所有节点畅通无阻地通信，就必须规定一个\uln{所有节点都遵守的协议}，也就是
TCP/IP。而银行内部的互连网不需要和外面互联互通，内部用什么协议由它自己决定。

\begin{bookbox}

王道书第 1 页补充：

\begin{itemize}
\tightlist
\item
  一句话区分：\uln{网络将多台计算机连接在一起，而互连网则将多个网络通过路由器连接在一起}；
\item
  \textbf{internet}（互连网）是\textbf{通用名词}，泛指由多个计算机网络互连而成的网络，可以用任意通信协议，并不要求必须采用
  TCP/IP；\textbf{Internet}（互联网或因特网）是\textbf{专用名词}，特指当前全球最大的、开放的、由众多网络和路由器互连而成的特定计算机网络，必须采用
  TCP/IP 协议族；
\item
  \textbf{ISP 的作用}：网络的最终用户通过 ISP
  接入互联网。所谓上网，是指通过某 ISP 分配的 IP 地址接入互联网。ISP
  从互联网管理机构申请大量 IP 地址（互联网上的主机都必须拥有 IP
  地址才能接入互联网，第 4
  章讨论），并拥有通信线路及路由器等联网设备。任何机构或个人只要向 ISP
  缴纳费用，就可获得所需 IP 地址的使用权，并通过该 ISP 接入互联网。
\end{itemize}

\end{bookbox}

\hypertarget{ux516dux56deux5230ux5f00ux5934ux7684ux56dbux4e2aux95eeux9898}{%
\section{六、回到开头的四个问题}\label{ux516dux56deux5230ux5f00ux5934ux7684ux56dbux4e2aux95eeux9898}}

\cueblock{\cuenote}{\relax{宽带断网、上网、投屏、银行内网，各指哪个“网”？}}

\slidefig{assets/1.1_1-p05.jpg}{0.5}{课件：以下提到的“网”指的是哪个网？}

{\rowcolors{2}{wdblue!5}{white}\arrayrulecolor{wdblue!35}
\begin{xltabular}{\linewidth}{>{\hsize=1.125\hsize\raggedright\arraybackslash}X>{\hsize=0.656\hsize\raggedright\arraybackslash}X>{\hsize=1.219\hsize\raggedright\arraybackslash}X}
\rowcolor{wdblue}\textcolor{white}{\bfseries 说法} & \textcolor{white}{\bfseries 指的是} & \textcolor{white}{\bfseries 依据} \\
\endfirsthead
① 宽带欠费，断网了 & 互联网 Internet & 欠费后 ISP 停止服务，家里接不上互联网（老师没有逐一给出这题的答案，这是按前面的讲解推出来的） \\ \hline
② 一起去上网 & 互联网 Internet & 老师原话：上网的“网”指互联网，也就是因特网 \\ \hline
③ 投屏要连同一个网络 & 计算机网络 & 手机和电视连在同一个（逻辑上的）交换机上 \\ \hline
④ 银行内网 & 互连网 internet & 各分行网络经路由器连接，不与外面互联互通 \\ \hline
\end{xltabular}\par\vspace{4pt}}

\hypertarget{ux4e03ux77e5ux8bc6ux56deux987e}{%
\section{七、知识回顾}\label{ux4e03ux77e5ux8bc6ux56deux987e}}

\slidefig{assets/1.1_1-p16.jpg}{0.6}{课件：知识回顾与重要考点}

\begin{summarybox}

\textbf{计算机网络}：由若干\textbf{节点}和连接这些节点的\textbf{链路}组成；节点可以是计算机、集线器、交换机、路由器等。主要功能：数据通信（信息传递）、资源共享。

\textbf{互连网（internet）}：多个计算机网络通过\hlred{路由器}互相连接而成；\hlred{可使用任意通信协议}。

\textbf{互联网（Internet，因特网）}：世界上最大的互连网，由各大 ISP
和国际机构组建，覆盖全球；\hlred{必须使用 TCP/IP 协议族}。

\textbf{其他}：ISP =
互联网服务提供商；集线器、交换机把多个节点连成同一个网络，路由器连接不同的网络；家用路由器
= 路由器 + 交换机 + 其他功能。

\textbf{自测}：① 定义里“自治”指什么？② internet 和 Internet
在协议上有什么区别？③ 交换机和路由器分别连接什么？④
为什么投屏要求手机和电视连同一个网络？

\end{summarybox}
\end{document}
